\documentclass[11pt]{article}
\usepackage{fontspec}
\usepackage[french,english]{babel}
\usepackage{amsmath}
\usepackage{amssymb}
\usepackage{amsthm}
\usepackage{mathrsfs}
\usepackage{tikz-cd}
\usepackage{enumitem}

\newcommand{\printedpage}[1]{\marginpar{\footnotesize p.~#1}}
\newcommand{\msp}[1]{\marginpar{\footnotesize [#1]}}
\newcommand{\uHom}{\underline{\mathrm{Hom}}}
\newcommand{\uFl}{\underline{\mathrm{Fl}}}
\newcommand{\uIs}{\underline{\mathrm{Is}}}
\newcommand{\uGamma}{\underline{\Gamma}}
\newcommand{\bDelta}{\boldsymbol{\Delta}}
\newcommand{\Ens}{\mathrm{Ens}}
\newcommand{\Ob}{\mathrm{Ob}\,}
\newcommand{\Fl}{\mathrm{Fl}}
\newcommand{\Hom}{\mathrm{Hom}}

\begin{document}

\printedpage{154}
% head: A. Grothendieck : Dérivateurs, Ch. XVIII, Édition des Universités de Montpellier II et Paris VII
in the category of double arrows of $\mathcal{E}_s$ ($I$ denoting the double-arrow category of type $0\rightrightarrows 1$), a functor which is (trivially) accessible. Let then $I\xrightarrow{\;i\;}\bDelta_1$ be the canonical functor, and
\[
\underbrace{\uHom(\bDelta_1,\mathcal{E}_s)}_{\simeq\,\uFl(\mathcal{E}_s)}\ \hookrightarrow\ \uHom(I,\mathcal{E}_s)
\]
the functor $u\longmapsto u\circ i$, which is a strictly fully faithful functor, having as its image the subcategory of double arrows $(u,v)$ such that $u=v$. As the inclusion functor is accessible and $\uHom(\bDelta_1,\mathcal{E}_s)$ and $\uHom(I,\mathcal{E}_s)$ accessible
\msp{page 406}
(a special case of 4.14.1, or also of 4.10.8!), we see that the strictly full subcategory of $\uHom(I,\mathcal{E}_s)$ of double arrows equalized in $\mathcal{E}_s$, is accessible in $\uHom(I,\mathcal{E}_s)$. So, its inverse image in $\mathcal{M}'$ is accessible. So the intersection of these inverse images for $s\in\Ob\mathcal{C}$ and $(\lambda,\mu)\in\Fl_2(\mathcal{C})$ varying, is so as well ($\mathcal{C}$ being small).

So everything comes down to proving that $\mathcal{M}'$ is accessible, which is contained in the

% vertical bar printed in the left margin alongside Lemme 4.14.22.1
\textbf{Lemma 4.14.22.1.} \textit{Let $\Pi$, $\Pi'$ be two accessible categories, $F,G:\Pi\rightrightarrows\Pi'$ two accessible functors, and $\mathcal{M}'$ the category of pairs $(X,u)$, with $X\in\Ob\Pi$ and $u:F(x)\longrightarrow G(x)$ an arrow of $\Pi'$. Then $\mathcal{M}'$ is accessible, and the canonical functor $\mathcal{M}'\longrightarrow\Pi$ is accessible.} % [sic] (X,u), X \in Ob Pi ... F(x) -> the same letter (X or x) throughout, capital X and lowercase x used for the same object

\textsc{Proof of the Lemma.} The datum of $\Pi$, $\Pi'$, $F$, $G$ defines a cofibered category $\mathcal{F}$ over the category $I=(0\rightrightarrows 1)$
\msp{page 407}
above. The category $\uGamma(\mathcal{F}/I)$ is interpreted, up to isomorphism, as the category of triples $(x,u,v)$ with $x\in\Ob\Pi$, $x'\in\Ob\Pi'$, and $u$ and $v$ arrows
\[
\begin{tikzcd}[column sep=small]
F(x) \arrow[dr, "u"'] & & G(x) \arrow[dl, "v"] \\
& x' &
\end{tikzcd}
\]
in $\Pi'$. The category $\mathcal{M}'$ is equivalent to the strictly full subcategory of $\uGamma(\mathcal{F}/I)$, consisting of triples $(x,u,v)$ such that $v$ is an isomorphism, i.e. to the subcategory of $\uGamma(\mathcal{F}/I)$, the inverse image of $\uIs(\Pi')$ under the functor $(x,u,v)\longmapsto v$,
\[
\uGamma(\mathcal{F}/I)\ \longrightarrow\ \uFl(\Pi')\ .
\]
Now by 4.10.8, $\uGamma(\mathcal{F}/I)$ is accessible, and the preceding functor is of course accessible. Finally, by 4.14.12,
\msp{page 408}
$\uIs(\Pi')$ is accessible in the accessible category $\uFl(\Pi')$. So its inverse image $\mathcal{M}'$ in $\uGamma(\mathcal{F}/I)$ is accessible in $\uGamma(\mathcal{F}/I)$, hence accessible. So the same holds for the equivalent category, q.e.d.

Thus, thanks to the general results on accessible parts of accessible categories, theorem 4.14.17, in the case where $p$ is fibrant, is reduced to the cofibrant cases, in the very special case where the index category is $I=(0\rightrightarrows 1)$. One could moreover make a similar reduction in the case where $p$ is cofibrant, by a proof practically % [sic] aux cas cofibrant -> au cas cofibrant, singular noun
% foot: Version 17 décembre 2025
\printedpage{155}
% head: A. Grothendieck : Dérivateurs, Ch. XVIII, Édition des Universités de Montpellier II et Paris VII
identical. In other words, theorem 4.10.8 reduces to the case where the index category is $I=(0\rightrightarrows 1)$. But I suspect that the proof of the theorem in this special case, is not essentially simpler than in the general case.

\msp{page 409}

\textbf{4.14.23.} One may ask about finding a statement ensuring the accessibility of a category $\uGamma(\mathcal{E}/\mathcal{C})$, which covers the two cases considered in 4.14.17. Here is what I can propose. I am going to make the following assumption on the functors $p:\mathcal{E}\longrightarrow\mathcal{C}$ (in addition to that of $\mathcal{C}$ being essentially small): % [sic] les foncteur -> le foncteur, singular functor p

\begin{itemize}[leftmargin=*]
\item[H)] For every $\lambda\in\Fl\mathcal{C}$, the category $\mathcal{E}_\lambda=\mathcal{E}\times_{\mathcal{C}}(\bDelta_1,i_\lambda)$, where $i_\lambda:\bDelta_1\longrightarrow\mathcal{C}$ is the functor defined by $\lambda$, is accessible, and the canonical functor $\mathcal{E}_\lambda\xrightarrow{\;p_\lambda\;}\bDelta_1$ is accessible.
\end{itemize}

Note that if $\lambda:s\longrightarrow t$ in $\mathcal{C}$ is such that $\mathcal{E}_\lambda$ and $p_\lambda$ are accessible, then $\mathcal{E}_s$ and $\mathcal{E}_t$ are accessible, since isomorphic to the inverse image categories under $p_\lambda:\mathcal{E}_\lambda\longrightarrow\bDelta_1$ of the subcategories $\{0\}$ resp. $\{1\}$ of $\Delta_1$, which are accessible in the accessible category $\bDelta_1$, so % [sic] sont accessible -> sont accessibles, plural agreement
\msp{page 410}
$\mathcal{E}_{\lambda_0}$ and $\mathcal{E}_{\lambda_1}$ are accessible in $\mathcal{E}_\lambda$, hence are accessible, hence also $\mathcal{E}_s$ and $\mathcal{E}_t$. Thus, assumption H) implies that the fiber categories are accessible. Note moreover:

\textbf{Lemma 4.14.23.1.} \textit{Suppose that $p:\mathcal{E}\longrightarrow\mathcal{C}$ is cofibrant or fibrant, and let $(\lambda:s\longrightarrow t)\in\Fl(\mathcal{C})$. For $\mathcal{E}_\lambda$ to be accessible and $\mathcal{E}_\lambda\xrightarrow{\;p_\lambda\;}\bDelta_1$ accessible, it is necessary and sufficient that $\mathcal{E}_s$, $\mathcal{E}_t$ be accessible and $\lambda_*:\mathcal{E}_s\longrightarrow\mathcal{E}_t$ resp. $\lambda^*:\mathcal{E}_t\longrightarrow\mathcal{E}_s$ be an accessible functor.}

\textsc{Proof.} We may suppose that $\mathcal{C}=\bDelta_1$, and that $\mathcal{E}$ is the cofibered (resp. fibered) category over $\bDelta_1$ defined by a functor $\varphi:\mathcal{E}_0\longrightarrow\mathcal{E}_1$ (resp. $\psi:\mathcal{E}_1\longrightarrow\mathcal{E}_0$), and to show that $\mathcal{E}$ and $p:\mathcal{E}\longrightarrow\bDelta_1$ are accessible if and only if $\mathcal{E}_0$ and $\mathcal{E}_1$ and $\varphi$ (resp. $\psi$) are.

\msp{page 411}

Let us prove the ``only if''. We have already seen that if $\mathcal{E}$, $p$ are accessible, $\mathcal{E}_0$ and $\mathcal{E}_1$ are accessible in $\mathcal{E}_s$, hence accessible, and the inclusion functors % [sic] dans E_s -> dans E, here C = Delta_1 and the subcategories lie in E
\[
\mathcal{E}_0\ \hookrightarrow\ \mathcal{E}\ \hookleftarrow\ \mathcal{E}_1
\]
accessible. To say that $p$ is cofibrant (resp. fibrant) means that the inclusion functor $\mathcal{E}_1\hookrightarrow\mathcal{E}$ admits a right adjoint $\bar{\varphi}$ (resp. that the functor $\mathcal{E}_0\hookrightarrow\mathcal{E}$ admits a left adjoint $\tilde{\psi}$). Then $\varphi$ (resp. $\psi$) is the composite $\mathcal{E}_0\longrightarrow\mathcal{E}\xrightarrow{\;\bar{\varphi}\;}\mathcal{E}_1$ (resp. the composite $\mathcal{E}_1\hookrightarrow\mathcal{E}\xrightarrow{\;\tilde{\psi}\;}\mathcal{E}_1$). So to prove that $\varphi$ (resp. $\psi$) is accessible, it suffices to prove that $\bar{\varphi}$ (resp. $\tilde{\psi}$) is. This is a special case of the % [sic] E -> E_1 (target of psi-tilde composite) -> E_0, psi goes from E_1 to E_0

\textbf{Lemma 4.14.23.2.} \textit{Let $\mathcal{E}$, $\mathcal{E}'$ be two accessible categories} (\footnote[172]{\textbf{N.B.} It suffices that $\mathcal{E}$ admit a small generating subcategory, and that the objects of $\mathcal{E}'$ be accessible.}), \textit{and $\mathcal{E}\underset{G}{\overset{F}{\rightleftarrows}}\mathcal{E}'$ a pair of adjoint functors between $\mathcal{E}$ and $\mathcal{E}'$. Then $F$ and $G$ are accessible.}

\msp{page 412}

Indeed, $F$ is accessible, since it commutes with arbitrary small $\varinjlim$, a fortiori it satisfies $\mathrm{L}_{\pi_0}$ (where $\pi_0$ is such that $\mathcal{E}$, $\mathcal{E}'$ satisfy $\mathrm{L}_{\pi_0}$). Let us prove that $G$ is accessible. If $S$ is a
% foot: Version 17 décembre 2025
\printedpage{156}
% head: A. Grothendieck : Dérivateurs, Ch. XVIII, Édition des Universités de Montpellier II et Paris VII
small generating set of elements in $\mathcal{E}$, so that the small family of functors
\[
h_\xi\ =\ (x\longmapsto\Hom_{\mathcal{E}}(\xi,x))\ :\ \mathcal{E}\ \longrightarrow\ \Ens
\]
is conservative, to prove that $G$ is accessible, it suffices to prove that the $h_\xi\circ G$ ($\xi\in S\subset\Ob\mathcal{E}$) are. Now
\[
L_\xi\circ G(x')\ =\ \Hom(\xi,G(x'))\ \simeq\ \Hom_{\mathcal{E}'}(F(\xi),x')\ , % [sic] L_xi -> h_xi, the functors were named h_xi
\]
and since $F(\xi)$ is an accessible object of $\mathcal{E}'$, the functor
\[
x'\ \longmapsto\ (L_\xi\circ G)(x')\ \simeq\ \Hom_{\mathcal{E}'}(F(\xi),x') % [sic] L_xi -> h_xi, the functors were named h_xi
\]
is indeed accessible, q.e.d.

This proves the ``only if'' in 4.14.23.1. Let us prove the ``if'', i.e. suppose $\mathcal{E}_s$, $\mathcal{E}_t$ accessible, and $\varphi$ (resp. $\psi$) accessible.

\msp{page 413}

In both cases, $\mathcal{E}_0$ is an open of $\mathcal{E}$, $\mathcal{E}_1$ is the complementary closed, and $\mathcal{E}$ is the category ``glued'' with the help of the bifunctor
\[
H\ :\ \mathcal{E}_0^{\circ}\times\mathcal{E}_1\ \longrightarrow\ \Ens
\]
given by
\[
\left\{
\begin{aligned}
H(x_0,x_1)\ &=\ \Hom(\varphi(x_0),x_1)\\
&\hspace{-2em}\text{resp.}\\
H(x_0,x_1)\ &=\ \Hom(x_0,\psi(x_1))\ ,
\end{aligned}
\right.
\]
according as we are in the case $p$ cofibrant or $p$ fibrant. We are going to use the

\textbf{Proposition 4.14.24} (\footnote[173]{cf. more general statements and nice proof, 4.15.1 and 4.15.4 (p. 425, 432).}). \textit{Let $\mathcal{E}$ be a category, $\mathcal{E}_0$ an open subcategory, $\mathcal{E}_1$ the complementary closed subcategory. Equivalent conditions:}
\begin{itemize}[leftmargin=*]
\item[\textit{a)}] \textit{$\mathcal{E}$ is accessible, and $\mathcal{E}\longrightarrow\bDelta_1$ (with fibers $\mathcal{E}_0$, $\mathcal{E}_1$) is accessible.}
\item[\textit{b)}] \textit{$\mathcal{E}$ is accessible, and $\mathcal{E}_0$, $\mathcal{E}_1$ are accessible in $\mathcal{E}$.}
\item[\textit{c)}] \textit{$\mathcal{E}_0$, $\mathcal{E}_1$ are accessible, and for every $x_0\in\Ob\mathcal{E}_0$, the functor}
\[
x_1\ \longmapsto\ \uHom_{\mathcal{E}}(x_0,x_1)\ :\ \mathcal{E}_1\ \longrightarrow\ \Ens
\]
\msp{page 414}
\textit{is accessible, and there exists $\pi$ such that for every $x_1\in\Ob\mathcal{E}_1$, the functor}
\[
x_0\ \longmapsto\ \uHom_{\mathcal{E}}(x_0,x_1)^{\circ}\ :\ \mathcal{E}_0\ \longrightarrow\ (\Ens)^{\circ}
\]
\textit{is $\pi$-accessible.}
\end{itemize}
% foot: Version 17 décembre 2025

\printedpage{157}
% head: A. Grothendieck : Dérivateurs, Ch. XVIII, Édition des Universités de Montpellier II et Paris VII
Let us admit this lemma \texttt{[rather this proposition]}, let us finish the proof of 4.14.23.1. We suppose $\mathcal{E}_0$, $\mathcal{E}_1$ and $\varphi$ (resp. $\psi$) accessible, we want to prove condition a) of 4.14.23.3. For this, it suffices to prove condition c). Now we have, in the case $p$ cofibrant, % [sic] 4.14.23.3 -> 4.14.24, the conditions a), c) are those of Proposition 4.14.24
\[
\Hom_{\mathcal{E}}(x_0,x_1)\ \simeq\ \Hom_{\mathcal{E}_1}(\varphi(x_0),x_1)\ ,
\]
this is indeed accessible in $x_1$, since $\mathcal{E}_1$ is accessible, and $\pi$-accessible in $x_0$ if $\varphi$ is, as a composite of $\varphi$ with $\mathcal{E}_1\longrightarrow\Ens$, $\Hom(-,x_1)$, which is. In the case $p$ fibrant, we have
\[
\Hom_{\mathcal{E}}(x_0,x_1)\ \simeq\ \Hom_{\mathcal{E}_0}(x_0,\psi(x_1))\ ,
\]
this is accessible in $x_1$, as the composite of the accessible functor $\psi:\mathcal{E}_1\longrightarrow\mathcal{E}_0$, and the likewise accessible functor $\Hom(x_0,-):\mathcal{E}_0\longrightarrow\Ens$, which is $\pi$-accessible in $x_0$ tautologically.

To finish the proof of 4.14.23.1, it therefore remains to prove proposition 4.14.24. We have already noted that we have a $\Longrightarrow$ b, and b $\Longrightarrow$ c clear, since the functor considered $\mathcal{E}\longrightarrow\Ens$ is the composite $\mathcal{E}_1\overset{\text{incl.}}{\hookrightarrow}\mathcal{E}\xrightarrow{\;\Hom_{\mathcal{E}}(x_0,-)\;}\Ens$, is accessible as a composite of two
\msp{page 415}
accessible functors. So it remains to prove c $\Longrightarrow$ a. Let then $\pi_0$ be a useful cardinal for $\mathcal{E}_0$, and such that the functors
\[
x_0\ \longmapsto\ H(x_0,x_1)\ :\ \mathcal{E}_0\longrightarrow\Ens^{\circ}
\]
are $\pi_0$-accessible, let $\pi$ be a cardinal $\geqslant\sup(\pi_0,\mathrm{card}\,\text{red}\,\mathcal{E}_{0\,\pi_0})$ such that $\pi^{\pi_0}=\pi$ (hence useful for $\mathcal{E}_0$), useful for $\mathcal{E}_1$, and such that the functors
\[
x_1\ \longmapsto\ H(x_0,x_1)\ :\ \mathcal{E}_1\longrightarrow\Ens
\]
are $\pi$-accessible, for $x_0\in\Ob\mathcal{E}_{0\,\pi_0}$. Let $\mathcal{E}(\pi)$ be the strictly full subcategory of $\mathcal{E}$, such that
\[
\Ob\mathcal{E}(\pi)\ =\ \Ob\mathcal{E}_{0\,\pi}\cup\Ob\mathcal{E}_{1\,\pi}\ .
\]
I say that
\begin{itemize}[leftmargin=*]
\item[1°)] $\mathcal{E}$ satisfies $\mathrm{L}_\pi$, and $\mathcal{E}_0$ and $\mathcal{E}_1$ are stable in $\mathcal{E}$ under $\mathrm{L}_\pi$, the functors $\mathcal{E}_0\longrightarrow\mathcal{E}$, $\mathcal{E}_1\longrightarrow\mathcal{E}$, $\mathcal{E}\longrightarrow\bDelta_1$ satisfy $\mathrm{L}_\pi$ (\footnote[174]{it suffices that $\mathcal{E}_0$ and $\mathcal{E}_1$ satisfy $\mathrm{L}_\pi$, and that
\[
x_0\ \longmapsto\ H(x_0,x_1)\ :\ \mathcal{E}_0\longrightarrow\Ens^{\circ}
\]
satisfy $\mathrm{L}_\pi$ $\forall\ x\in\Ob\mathcal{E}$.}).
\item[2°)] $\mathcal{E}(\pi)=\mathcal{E}_\pi$ (whence a fully faithful functor
\[
\mathrm{Ind}_\pi(\mathcal{E}(\pi))\ \longrightarrow\ \mathcal{E}\ \ ).
\]
\item[3°)] The said functor is essentially surjective,
\msp{page 416}
hence an equivalence of categories.
\end{itemize}
Note that 2°) and 3°) imply that $\mathcal{E}$ is accessible, and 1°) that $\mathcal{E}\longrightarrow\bDelta_1$ is accessible, whence condition a) of the lemma we are to establish. So everything comes down to proving 1°, 2°, 3° above.
% foot: Version 17 décembre 2025

\printedpage{158}
% head: A. Grothendieck : Dérivateurs, Ch. XVIII, Édition des Universités de Montpellier II et Paris VII
1°) Since $\mathcal{E}_1$ is closed in $\mathcal{E}$, it is clear that $\mathcal{E}$ is stable under all $\varinjlim$ representable in $\mathcal{E}$, and more precisely, that for a functor $I\longrightarrow\mathcal{E}_1$, its direct limit in $\mathcal{E}_1$ exists if and only if it exists in $\mathcal{E}$, and the two are isomorphic. Suppose $I$ large relative to $\pi$, and let $\varphi:I\longrightarrow\mathcal{E}$, $i\longmapsto x_i$. Let us distinguish two cases.

a) $\varphi$ takes its values in $\mathcal{E}_0$. Let $x_0$ be its direct limit \textit{in} $\mathcal{E}_0$. I say that this is also a direct limit in $\mathcal{E}$. We must prove that for $x_1\in\Ob\mathcal{E}_1$, the arrow
\[
\Hom_{\mathcal{E}}(x_0,x_1)\ \longrightarrow\ \varprojlim_{i\in I}\Hom_{\mathcal{E}}(x_i,x_1)
\]
is isomorphic. Now this is the case
\msp{page 417}
by the assumption on $\pi$,
\[
\xi\ \longmapsto H(\xi,x_1)\ :\ \mathcal{E}_0\ \longrightarrow\ (\Ens)^{\circ}
\]
satisfies $\mathrm{L}_\pi$ for every $x_1$ in $\Ob\mathcal{E}_1$.

b) $\varphi$ does not factor through $\mathcal{E}_0$, i.e. $\exists\ i_0\in\Ob I$ with $\varphi(i_0)\in\Ob\mathcal{E}_1$. Then we have $\varphi(i)\in\Ob\mathcal{E}_1$ for every $i\in\Ob I[i_0]$, where $I[i_0]$ denotes the strictly full subcategory of $I$ consisting of the $i$ such that $i\geqslant i_0$. This category is cofinal in $I$, hence large relative to $\pi$ like $I$, and for $\varinjlim\limits_{I}x_i$ to exist, it is necessary and sufficient that $\varinjlim\limits_{I[i_0]}x_i$ exist, which is the case by the cases a) already treated. % [sic] les cas a) déjà traité -> le cas a) déjà traité, singular agreement

This therefore proves at the same time that $\mathcal{E}$ satisfies $\mathrm{L}_\pi$, and that $\mathcal{E}_0$ and $\mathcal{E}_1$ are $\mathrm{L}_\pi$-stable. The proof shows at the same time that the functor $p:\mathcal{E}\longrightarrow\bDelta_1$ satisfies $\mathrm{L}_\pi$.

2°) From this it follows at once that $\mathcal{E}_\pi\subset\mathcal{E}(\pi)$. Let us prove that we also have
\[
\mathcal{E}(\pi)\ \subset\ \mathcal{E}_\pi\ ,
\]
and let us again distinguish two cases, for
\msp{page 418}
$x\in\Ob E(\pi)$, according as $x$ is in $\mathcal{E}_{0\,\pi}$ or $\mathcal{E}_{1\,\pi}$. If it is in $\mathcal{E}_{1\,\pi}$, to prove that % [sic] E(pi) -> \mathcal{E}(pi), script letter elsewhere
\[
\Hom_{\mathcal{E}}(x,\underbrace{\varinjlim y_i}_{y})\ \xleftarrow{\ \sim\ }\ \varinjlim_{i}\Hom_{\mathcal{E}}(x,y_i)
\tag*{$(*)$}
\]
(for a direct system $(y_i)_{i\in I}$ in $\mathcal{E}$ large relative to $\pi$), if the $y_i$ are in $\mathcal{E}_0$, hence $y\in\mathcal{E}_0$, then both sides are empty, and this is OK. Otherwise, we may suppose (replacing $I$ by an $I[i_0]$ if need be) that the $x_i$ are in $\mathcal{E}_1$, and then this is OK too, since $y$ is also the $\varinjlim$ in $\mathcal{E}_1$, by 1°). % [sic] les x_i -> les y_i, the inductive system is (y_i)

If $x$ is in $\mathcal{E}_{0\,\pi}$, then $(*)$ is clear if the $y_i$ are in $\mathcal{E}_0$. Otherwise, we may again suppose that the $y_i$ are in $\mathcal{E}_0$, and $(*)$ reads % [sic] y_i sont dans E_0 (second occurrence) -> E_1, the case "sinon" and the following "les y_i dans E_1"
\[
H(x,\varinjlim y_i)\ \xleftarrow{\ \sim\ }\ \varinjlim H(x,y_i)\ ,
\]
for $x\in\Ob\mathcal{E}_{0\,\pi}$, the $y_i$ in $\mathcal{E}_1$, $I$ large relative to $\pi$. In other words, we must prove that the functor
\[
x_1\ \longmapsto\ H(x_0,x_1)\ :\ \mathcal{E}_1\ \longrightarrow\ \Ens
\]
% foot: Version 17 décembre 2025

\end{document}
